\section{Introduction}
\label{sec:intro}

The automated monitoring of agricultural infrastructure from high-resolution aerial imagery is an urgent imperative for global food security, environmental policy enforcement, and epidemiological surveillance. In particular, industrial poultry farming---characterized by geographically dispersed and geometrically specific structural footprints---requires highly scalable mapping systems capable of functioning at continental scales. While supervised deep learning frameworks have achieved significant success in remote sensing object detection and segmentation \cite{li2020greenhouse, robinson2021mapping}, their generalization is inherently bottlenecked by the immense cost and labor required to acquire dense annotations across heterogeneous geographic regions and temporal shifts.

Traditionally, remote sensing models are initialized with weights pre-trained on ImageNet. However, the pronounced domain gap between ground-level natural photographs and top-down aerial imagery severely limits the efficacy of these generic priors \cite{wang2022empirical}. Aerial scenes exhibit arbitrary object orientations, extreme scale variance, and complex multispectral statistics that are fundamentally absent in natural-image benchmarks. To resolve this source-target mismatch, the Earth Observation (EO) community has increasingly embraced domain-specific foundation models, seeking to learn intrinsic semantic structures directly from massive, unlabeled remote sensing corpora \cite{sun2023ringmo}.

Currently, Masked Image Modeling (MIM) stands as the dominant paradigm for pre-training these foundation models \cite{tian2023spark, reed2022scale}. Yet, despite scaling efficiently to massive unannotated datasets, vanilla MIM faces severe limitations in applied detection contexts. By reconstructing highly masked images patch-by-patch, MIM inherently prioritizes local pixel fidelity, often smoothing over the distinct structural boundaries and small objects that are prevalent in dense aerial scenes. Conversely, self-distillation learning paradigms maintain highly separable global semantics \cite{caron2021emerging, jain2022multimodal}, proving exceptionally robust for dense structural mapping.

In this paper, we propose a domain-adaptive self-supervised learning pipeline tailored for agricultural infrastructure detection. By pre-training a Vision Transformer (ViT-Base/16) on roughly 9.87 million unannotated NAIP (National Agriculture Imagery Program) tiles using the DINO (self-distillation with no labels) framework \cite{caron2021emerging}, we engineer a robust geospatial foundation model. Our approach exclusively utilizes DINO's cross-entropy objective and multi-crop spatial augmentations to distill robust global semantics directly from the Earth Observation data distribution without relying on patch reconstruction.

Our primary contributions are fourfold:
\begin{enumerate}
    \item We present a high-throughput self-distillation pre-training pipeline for Earth Observation, processing a massive 1.5 TB NAIP corpus efficiently via streaming WebDatasets.
    \item We establish that our pure DINO self-supervised ViT-Base backbone averts dimensional collapse and maintains robust global semantics.
    \item We establish that our self-supervised ViT-Base backbone achieves highly competitive end-to-end downstream performance on poultry farm segmentation, reaching a robust mAP@0.5 of $0.977$ and matching heavily optimized ImageNet baselines.
    \item We rigorously ablate the sample efficiency of our representations in heavily constrained few-shot scenarios (10\% labeled data), exposing critical limitations regarding the label dependency of high-capacity self-distillation architectures under extreme data scarcity.
\end{enumerate}
